\section{Introduction}

In one of the games we study, only two players remain. Gemini writes in its private note that Sol's holdings can raise \$1,414 if everything is sold and mortgaged, then builds a hotel that lifts one rent to \$1,400. Sol lands there with \$49. Its note reads \Q{g20-plan} Sol sells its buildings, mortgages eight deeds and pays. It has \$14 left. In another game Sol owes \$500, adds up \$408 and writes \Q{g12-false-note} It declares bankruptcy. Its holdings could have raised \$748.

Both responses answer one question: what can be raised? An agent that manages money must often value resources that are not on the screen in front of it, such as collateral that an earlier action must free or the cash a counterparty will hold after a sale. Monopoly makes such valuations checkable to the dollar, because the board fixes every mortgage and resale value and the engine records every state.

This paper asks what four configured LLM systems count as raisable money in twenty complete four-player games: their own when a debt arrives and a rival's when a trade is proposed. We check each count against board arithmetic and compare it with the action the engine applied. The populations are complete. They are all 188 debts that required financing and all 920 responses to trade offers.

\textbf{The environment does the route finding.} Agents paid 128 of 131 payable debts, including fourteen of sixteen that needed mortgages released by building sales. We show from the engine source and the recorded holdings that in all 131 entry states, any uninterrupted sequence of valid cash-raising actions pays (\S\ref{sec:task}). Payment therefore records that an agent did not stop. The informative records are the first choice to continue or stop, the text written with it, and the assets given up.

\textbf{First assessments are often exact.} The financing menu lists deeds without prices and the prompt states no rules. On unpayable debts, 43 of 48 first responses state the full asset bound, its proceeds or the remaining gap to the dollar. In the sixteen payable debts that need released collateral, five responses state the dependency, three give an ordered route, five give a budget that includes it and one raises it as a question (\S\ref{sec:texts}). Agents also choose what to give up. When listed mortgages were enough, they kept every building in 28 of 30 payments.

\textbf{The exceptions have a specific form.} Five first responses, from all four systems, stop at the total the first menu lists. Three concern debts that were unpayable either way. Two concern payable debts, and both agents declared. Of nine first actions that depart from the asset bound, six are Grok tool calls that contradict the action named in the same response, and one is a mortgage that Gemini justifies by a published settlement rule the engine does not implement (\S\ref{sec:stop}).

\textbf{Agents apply the same accounting to rivals.} In bargaining they price a deed by what it lets the buyer build, refuse their own monopoly because of the other side's budget, and finance a rival so that its estate does not pass to a stronger creditor (\S\ref{sec:bargain}). This evidence is exploratory.

The study is observational. It describes one roster under unequal response conditions, with no intervention, no executed reference policy and no human annotation. It does not rank models, identify why an assessment was incomplete, or measure whether a bargaining tactic worked. Appendix~\ref{app:scope} maps the project's original research questions to what this corpus does and does not establish.
